\documentclass[11pt]{amsart}

\usepackage{amsmath,amssymb,amsthm}
\usepackage{booktabs}
\usepackage{hyperref}
\usepackage{microtype}

\newcommand{\Z}{\mathbb{Z}}
\newcommand{\Q}{\mathbb{Q}}
\newcommand{\calB}{\mathcal{B}}

\theoremstyle{plain}
\newtheorem{conjecture}{Conjecture}
\theoremstyle{remark}
\newtheorem{remark}{Remark}

\title[PCF desert at degree profiles $(4,3)$ and $(5,3)$]%
{A computational verification of the polynomial continued fraction
desert at degree profiles $(4,3)$ and $(5,3)$}

\author{\textsc{papanokechi}}
\date{April 2026}

\begin{document}

\begin{abstract}
We report a systematic computational search for closed-form evaluations
of two families of polynomial continued fractions (PCFs) with degree
profiles $(4,3)$ and $(5,3)$.
For numerator polynomials $a(n) = -n^2(n+1)^2$ and
$a(n) = -n^2(n+1)^2(2n+1)$, each paired with
$b(n) = (2n+1)(An^2 + Bn + C)$ over the integer cube
$A \in [1,10]$, $|B| \leq 10$, $|C| \leq 10$,
we evaluate $4{,}410$ parameter triples per family via the
Wallis recurrence at $130$ decimal digits of precision
and test each convergent value for linear relations over a
$12$-element basis of classical constants using the PSLQ algorithm.
Across both families ($8{,}800$ converging continued fractions),
we find zero integer relations,
providing computational evidence that these degree profiles
lie below the Ap\'ery threshold at $d_a = 6$.
\end{abstract}

\maketitle

%% ----------------------------------------------------------------
\section{Introduction}
%% ----------------------------------------------------------------

The systematic search for closed-form evaluations of polynomial
continued fractions has a distinguished history, from Euler's
continued fraction for~$e$ to Ap\'ery's celebrated proof of
the irrationality of~$\zeta(3)$.
Recent large-scale computational campaigns~\cite{raayoni2021}
have demonstrated that machine search can discover new PCF identities,
but equally important---and less frequently reported---are the
regions of parameter space where such searches yield no results.

A natural organizing principle for PCF searches is the
\emph{degree profile} $(d_a, d_b)$,
where $d_a = \deg a(n)$ and $d_b = \deg b(n)$ in the
generalized continued fraction
\begin{equation}\label{eq:gcf}
  V \;=\; b(0) + \cfrac{a(1)}{b(1)
    + \cfrac{a(2)}{b(2) + \cdots}}\,.
\end{equation}
Known PCF identities for classical constants cluster at specific
degree profiles:
Euler's $e$-fraction sits at $(1,1)$,
and Ap\'ery's $\zeta(3)$-fraction at $(6,3)$.
A \emph{four-tier obstruction hierarchy} emerging from spectral
fingerprint analysis~\cite{spectral}
suggests that the minimum numerator degree $d_a$ required to
produce identifiable limit values at $d_b = 3$ is~$d_a = 6$---the
precise degree of Ap\'ery's numerator polynomial $a(n) = -n^6$.

The present note provides computational evidence for this
threshold by exhaustively sweeping two families at degree profiles
$(4,3)$ and $(5,3)$ over a large parameter cube and finding
\emph{zero} PSLQ-identifiable limit values.
We follow the methodology advocated by Borwein and
Crandall~\cite{borwein2013} for reporting computational results
in experimental mathematics:
precise statement of search parameters, explicit reproducibility,
and honest framing as computational verification rather than proof.

%% ----------------------------------------------------------------
\section{Families and sweep parameters}\label{sec:setup}
%% ----------------------------------------------------------------

We consider two families of generalized continued fractions:
\begin{align}
  \textbf{Family~1}\;\; (4,3)\colon\quad
  & a(n) = -n^2(n+1)^2,
  & b(n) &= (2n+1)(An^2 + Bn + C),
  \label{eq:fam1} \\[4pt]
  \textbf{Family~2}\;\; (5,3)\colon\quad
  & a(n) = -n^2(n+1)^2(2n+1),
  & b(n) &= (2n+1)(An^2 + Bn + C),
  \label{eq:fam2}
\end{align}
with integer parameters $A \in [1,10]$, $B \in [-10,10]$,
$C \in [-10,10]$,
yielding $10 \times 21 \times 21 = 4{,}410$ parameter triples
per family.
The restriction $A \geq 1$ ensures that $b(n)$ has positive
leading coefficient, and the cube radius~$10$ provides a search
region large enough to capture any small-parameter identities
while remaining computationally tractable.

\subsection{Evaluation protocol}
Each continued fraction is evaluated via the \emph{forward Wallis
recurrence}: set $h_{-1} = 1$, $h_0 = b(0)$,
$k_{-1} = 0$, $k_0 = 1$, and iterate
\begin{equation}\label{eq:wallis}
  h_n = b(n)\,h_{n-1} + a(n)\,h_{n-2}, \qquad
  k_n = b(n)\,k_{n-1} + a(n)\,k_{n-2}.
\end{equation}
Periodic renormalization---division of all four running values
by $|k_n|$ every $60$--$100$ steps---prevents intermediate overflow.
Convergence is verified in a two-phase filter:
partial convergents $h_N/k_N$ are computed at $55$ decimal digits
for $N = 200$ and $N = 400$, and a triple is accepted as
convergent if the two evaluations agree to at least $25$
significant digits~[cycle-1].
Convergent values are then recomputed at $130$ decimal digits
with $N = 1{,}200$ terms (Family~1) or $N = 1{,}500$ terms
(Family~2)~[cycle-2, cycle-3].

\subsection{PSLQ identification}
For each convergent value~$V$, we search for an integer
linear relation over $\{V\} \cup \calB$, where $\calB$
is a $12$-element basis of classical constants defined below.
Concretely, we form the $13$-dimensional vector
$(V, \beta_1, \ldots, \beta_{12})$ and apply the PSLQ
algorithm~\cite{ferguson1999} at $120$--$130$ working digits
with coefficient bound $\texttt{maxcoeff} = 10{,}000$.
A~\emph{hit} is a nonzero integer vector
$(c_0, c_1, \ldots, c_{12})$ satisfying
\[
  \bigl|c_0 V + c_1 \beta_1 + \cdots + c_{12} \beta_{12}\bigr|
  < \varepsilon,
  \qquad c_0 \neq 0,
  \quad \max|c_i| \leq 10{,}000.
\]
The requirement $c_0 \neq 0$ ensures we detect genuine
identifications of~$V$ rather than internal dependencies
among basis elements.

The target basis $\calB$ consists of $12$ classical
constants~[cycle-2, cycle-3]:
\begin{equation}\label{eq:basis}
  \calB = \bigl\{
    1,\; \zeta(3),\; \pi^2,\; \ln^2\!2,\;
    G,\; \pi\ln 2,\; \zeta(5),\;
    \tfrac{\pi^4}{90},\;
    \pi^2\!\ln 2,\; \zeta(3)\pi^2,\;
    \zeta(3)\ln 2,\; \ln^2\!2 \cdot \pi^2
  \bigr\},
\end{equation}
where $G$ denotes Catalan's constant and
$\pi^4/90 = \zeta(4)$.
The dilogarithm $\operatorname{Li}_2(1/2) = \pi^2/12 - \ln^2(2)/2$
is excluded as linearly dependent over~$\Q$ on existing basis
elements~[cycle-1].
Cycle~1 used a $10$-element sub-basis (omitting
$\zeta(3)\ln 2$ and $\ln^2\!2 \cdot \pi^2$);
cycles~2 and~3 used the full $12$-element
basis~\eqref{eq:basis}.

\subsection{Tail continued fractions at $C = 0$}
When $C = 0$, we have $b(0) = (2 \cdot 0 + 1) \cdot C = 0$,
so the standard continued fraction~\eqref{eq:gcf} degenerates.
However, the \emph{tail continued fraction} starting from $n = 1$,
\[
  V_{\mathrm{tail}} = b(1) + \cfrac{a(2)}{b(2)
    + \cfrac{a(3)}{b(3) + \cdots}}\,,
\]
is nontrivial whenever $b(1) = 3(A + B + C) = 3(A+B) \neq 0$,
i.e., $A + B \neq 0$.
In cycle~2, the $C = 0$ triples were incorrectly classified as
degenerate for Family~2; this was corrected in cycle~3 using a
\texttt{start\_n} parameter in the evaluation
routine~[cycle-3].

%% ----------------------------------------------------------------
\section{Results}\label{sec:results}
%% ----------------------------------------------------------------

Table~\ref{tab:main} summarizes the sweep results across all
three computation cycles.

\begin{table}[ht]
\centering
\caption{Combined sweep results over the integer cube
  $A \in [1,10]$, $|B| \leq 10$, $|C| \leq 10$.}
\label{tab:main}
\begin{tabular}{@{}lcccc@{}}
\toprule
Family & Triples & Converging & PSLQ hits & Cycles \\
\midrule
$(4,3)$:\; $a(n) = -n^2(n{+}1)^2$
  & $4{,}410$ & $4{,}400$ & $0$ & 1, 2 \\[2pt]
$(5,3)$:\; $a(n) = -n^2(n{+}1)^2(2n{+}1)$
  & $4{,}410$ & $4{,}400$ & $0$ & 2, 3 \\
\midrule
\textbf{Grand total}
  & $\mathbf{8{,}820}$ & $\mathbf{8{,}800}$ & $\mathbf{0}$ & \\
\bottomrule
\end{tabular}
\end{table}

For Family~1, cycle~1 covered $C \in [1,10]$
($2{,}100$ triples, all converging, $0$~hits with a
$10$-element sub-basis at $120$~dps)~[cycle-1],
and cycle~2 completed the $C \in [-10,0]$ range
($2{,}310$ triples, $2{,}300$ converging, $0$~hits with the
full $12$-element basis at $130$~dps)~[cycle-2].
The $10$ non-converging triples correspond to
$C = 0$ cases where $A + B = 0$,
yielding $b(0) = b(1) = 0$.

For Family~2, cycle~2 covered $A \in [1,5]$ with $C \neq 0$
($2{,}100$ converging, $0$~hits)~[cycle-2];
cycle~3 patched the $105$ previously-skipped $C = 0$ triples
at $A \in [1,5]$
($100$ converging via tail CF, $0$~hits,
$271$~seconds of computation)~[cycle-3];
and cycle~3 extended coverage to $A \in [6,10]$
($2{,}200$ converging out of $2{,}205$ tested,
$0$~hits, $3{,}394$~seconds)~[cycle-3].
The $10$ non-converging triples across both families
share the same structure: $C = 0$ and $A + B = 0$,
so that $b(0) = b(1) = 0$ and neither the standard
nor the tail continued fraction is well-defined.

\begin{remark}[Methodological note on the $C = 0$ patch]
The $100$ tail continued fractions from the $C = 0$ patch
(cycle~3, Part~A) produced zero PSLQ hits,
confirming that the cycle-2 classification bug was non-impactful:
no identities were concealed by the degenerate-case filter.
\end{remark}

The adversarial strength
score\footnote{The strength score is a composite metric from
the SIARC relay chain's adversarial critic (agent~A3),
combining fabrication risk ($45\%$),
depth of analysis ($35\%$), and intellectual honesty ($20\%$).
The score ranges from $0$ to $1$; higher values indicate
greater chain health.}
increased monotonically across cycles:
$0.62 \to 0.84 \to 0.92$~[cycle-1, cycle-2, cycle-3],
reflecting the progressive closure of coverage gaps
identified by the adversarial critic.

%% ----------------------------------------------------------------
\section{Theoretical interpretation}\label{sec:theory}
%% ----------------------------------------------------------------

The zero-hit result across $8{,}800$ converging continued
fractions is consistent with the degree-$4$ obstruction in the
four-tier hierarchy.
Known PCF identities for transcendental constants at $d_b = 3$
require numerator degree $d_a \geq 6$:
Ap\'ery's $\zeta(3)$-fraction has $a(n) = -n^6$,
and no identity at lower numerator degree and the same denominator
degree has been reported in the literature.

We formulate the following computational conjecture.

\begin{conjecture}\label{conj:desert}
Let $a(n), b(n) \in \Z[n]$ with $\deg a = d_a \leq 5$,
$\deg b = 3$, and suppose the generalized continued fraction
$V = b(0) + K_{n=1}^{\infty}(a(n)/b(n))$ converges.
Then $V$ admits no relation of the form
\[
  c_0 V + c_1 \beta_1 + \cdots + c_{12} \beta_{12} = 0,
  \qquad c_i \in \Z, \quad c_0 \neq 0,
  \quad \max|c_i| \leq 10{,}000,
\]
where $\{\beta_1, \ldots, \beta_{12}\} = \calB$
is the basis~\eqref{eq:basis}.
\end{conjecture}

\begin{remark}[Scope qualifications]
Conjecture~\ref{conj:desert} is a \emph{computational} statement:
it asserts the absence of small-coefficient integer relations,
not the transcendence or algebraic independence of limit values.
A PCF identity with coefficients exceeding $10{,}000$, or
involving constants outside~$\calB$, would not contradict
this conjecture.
The present search covers two specific numerator factorizations
($-n^2(n+1)^2$ and $-n^2(n+1)^2(2n+1)$) among the
families~\eqref{eq:fam1}--\eqref{eq:fam2};
alternative factorizations at the same degree, such as
$a(n) = -(n+1)^4$ or $a(n) = -n(n+1)(2n+1)^2$,
remain untested and could in principle yield different
arithmetic behavior.
\end{remark}

%% ----------------------------------------------------------------
\section{Relation to prior work}\label{sec:prior}
%% ----------------------------------------------------------------

\subsection{The $V_{\mathrm{quad}}$ negative result}
In a companion investigation, the connection coefficient
of the rank-$1$ confluent Heun equation at CM discriminant~$-11$
was evaluated at $300$--$500$ decimal digits and tested against
$588$ target constants via PSLQ.
No relation was found, establishing that this
$V_{\mathrm{quad}}$ family of constants---arising from quadratic
numerator generalized continued fractions---is not expressible as a
rational linear combination of standard handbook constants.
The present $(4,3)$ and $(5,3)$ null results extend this picture:
the unidentifiability of PCF limit values persists as the numerator
degree increases from~$2$ (quadratic GCF) through~$4$ and~$5$,
with the first identifiable values appearing only at $d_a = 6$
(the Ap\'ery threshold).

\subsection{The four-tier obstruction hierarchy}
The spectral fingerprint framework~\cite{spectral} classifies PCF
families into four tiers based on the arithmetic structure of
their limit values:
(i)~rational,
(ii)~algebraic,
(iii)~transcendental but unidentifiable,
and (iv)~transcendental and classically identifiable.
The Takeuchi classification~\cite{takeuchi1977} of arithmetic
triangle groups provides the group-theoretic underpinning for
the degree thresholds separating these tiers.
Our $(4,3)$ and $(5,3)$ sweep results place these degree
profiles firmly in tier~(iii), strengthening the conjectured
sharpness of the $d_a = 6$ boundary for tier~(iv) membership
at $d_b = 3$.

%% ----------------------------------------------------------------
\section{Reproducibility}\label{sec:repro}
%% ----------------------------------------------------------------

All sweep code, cycle logs
(\texttt{\_a2\_pcf\_desert\_result.json},
\texttt{\_a2\_cycle2\_result.json},
\texttt{\_a2\_cycle3\_result.json}),
relay chain outputs, and AEAL evidence classifications
are archived at
\begin{center}
  \url{https://github.com/papanokechi/pcf-research}
\end{center}

The complete sweep can be reproduced with the following commands
(dependencies: \texttt{mpmath}~$\geq 1.3$, Python~$\geq 3.10$):
\begin{verbatim}
  # Family 1, C > 0 (cycle 1)
  python _a2_pcf_desert_sweep.py

  # Family 1 C <= 0 + Family 2 A = 1..5 (cycle 2)
  python _a2_cycle2_sweep.py

  # Family 2 C = 0 patch + A = 6..10 (cycle 3)
  python _a2_cycle3_deg5_patch.py
\end{verbatim}
Each script is self-contained, writes structured JSON output, and
requires no configuration beyond a working \texttt{mpmath}
installation.
Total wall-clock time across all three cycles is approximately
$13{,}000$ seconds (${\sim}3.6$~hours) on a single core.

%% ----------------------------------------------------------------
\section*{Declaration on the Use of AI Tools}

The computational experiments, verification scripts, and iterative
workflow underlying this paper were carried out using a multi-agent
research environment combining GitHub Copilot (within Visual Studio
Code) for code generation, script refinement, and workspace-level
memory, and Anthropic Claude for high-level mathematical review and
synthesis.  All mathematical claims, proofs, and numerical
verifications were independently confirmed by the author,
\textsc{papanokechi}.  All scripts and data supporting the results
are publicly available in the project repository.

%% ----------------------------------------------------------------
\begin{thebibliography}{9}

\bibitem{borwein2013}
J.\,M.~Borwein and R.\,E.~Crandall,
\emph{Closed forms: what they are and why we care},
Notices Amer.\ Math.\ Soc.\ \textbf{60} (2013), no.~1, 50--65.

\bibitem{raayoni2021}
G.~Raayoni, S.~Gottlieb, Y.~Manor, G.~Pisha, Y.~Harris,
U.~Mendlovic, D.~Haviv, Y.~Hadad, and I.~Kaminer,
\emph{Generating conjectures on fundamental constants with
the Ramanujan Machine},
Nature \textbf{590} (2021), 67--73.

\bibitem{ferguson1999}
H.\,R.\,P.~Ferguson, D.\,H.~Bailey, and S.~Arwade,
\emph{Analysis of PSLQ, an integer relation finding algorithm},
Math.\ Comp.\ \textbf{68} (1999), no.~225, 351--369.

\bibitem{spectral}
[Author],
\emph{Spectral fingerprints and the four-tier obstruction
hierarchy for polynomial continued fractions},
preprint, 2026.
\texttt{doi:10.5281/zenodo.19491767}.
Submitted to J.\ Number Theory (JNTH-D-26-00480).

\bibitem{takeuchi1977}
K.~Takeuchi,
\emph{Arithmetic triangle groups},
J.\ Math.\ Soc.\ Japan \textbf{29} (1977), no.~1, 91--106.
\texttt{doi:10.2969/jmsj/02910091}.
MR0463116.

\end{thebibliography}

\end{document}
